\cleardoublepage
\chapter*{Preface}
\markboth{Preface}{Preface}
Biostatistics is a work in progress. It is based on a collection of expanding lecture notes on statistics as used by biologists such as botanists, zoologists, and ecologists. It is intended to be a resource for students and researchers who are interested in learning more about the subject and applying it to their data. The material is presented in a way that is accessible to those with a thorough understanding of biology and a willingness and interest to learn the application of statistical analyses to their data.

Although the book is primarily aimed at biologists, it is also suitable for individuals in other fields who require subjecting their data to inferential statistical analyses. The material is presented in a way that is accessible even to those who might not have a deep interest of background in mathematics or statistics.

The examples in this book require that the reader already possesses a good working knowledge of the R programming language. If you are not familiar with R, I recommend that you take a look at the excellent book \href{https://r4ds.had.co.nz/}{R for Data Science} by Hadley Wickham and Garrett Grolemund. You may also study the material provided on \href{https://tangledbank.netlify.app}{The Tangled Bank} website.

You are also required to already have a good grasp of exploratory data analysis (EDA), descriptive (summary) statistics, and data visualisation. If you are not familiar with these topics, I recommend that you take a look at the excellent book \href{https://bookdown.org/rdpeng/exdata/}{Exploratory Data Analysis with R} by Roger D. Peng. \href{https://tangledbank.netlify.app}{The Tangled Bank} is also there to help you.

The book is divided into several chapters, each of which covers a different aspect of statistics. The chapters are organised in a logical order, starting with the basics of probability and moving on to more advanced topics such as hypothesis testing, regression analysis, and multivariate statistics (eventually, maybe, and possible, but there are already excellent text available on the topic so it is not a priority for me to write this material yet). Each chapter is accompanied by a set of exercises that are designed to help the reader understand the material and apply it to real-world problems.

Biostatistics shows you to connect i) the biological question, ii) the data you collect, and iii) the model you use to analyse it. When those three components fit together, your analysis becomes coherent and defensible. Throughout the book, therefore, you will see how to move from a question about a biological system to a structured analysis. This includes summarising data, selecting appropriate methods, fitting models, checking whether those models are adequate, and communicating your findings to a scientific audience.